% !TeX root = ../navier-stokes-manuscript.tex
% ============================================================
% 1.6 A Comparison with Euler
% ============================================================
\subsection{A Comparison with Euler}\label{sub:AComparisonWithEuler}

The Couette picture shows tangential momentum passing through neighbouring layers even when the net viscous force inside each layer is zero.
Remove viscosity from the same kind of tangential scene and a sharp slip becomes possible.

Two rigid boards sliding along a common face give the kinematic picture.
Their normal velocities agree because neither board passes through the other, while their tangential velocities can differ abruptly at the interface.
The fluid analogue is an idealized vortex sheet.

Let
\[
  \Sigma=\{x_3=0\},
\]
choose constants \(U_-\neq U_+\), and define
\[
  u(x)=
  \begin{cases}
    U_-e_1, & x_3<0,\\
    U_+e_1, & x_3>0,
  \end{cases}
  \qquad
  p(x)=p_0.
\]
The jump
\[
  [u]=(U_+-U_-)e_1
\]
is tangent to \(\Sigma\), whose unit normal is \(e_3\).
The distributional divergence contains the possible sheet term \([u]\cdot e_3\,\delta_\Sigma\), and that term vanishes because
\[
  [u]\cdot e_3=0.
\]

The stationary Euler momentum equation is just as direct in divergence form:
\[
  \nabla\cdot(u\otimes u)+\nabla p=0.
\]
On both sides of the sheet,
\[
  (u\otimes u)e_3=u(u\cdot e_3)=0.
\]
There is no jump in normal momentum flux, and the constant pressure contributes no sheet force.
The piecewise constant field is a stationary weak solution of the incompressible Euler equations.

Its loss of smoothness is concentrated at the interface.
Since \(u_1\) jumps in the \(x_3\)-direction,
\[
  \nabla\times u
  =
  (U_+-U_-)e_2\,\delta_\Sigma.
\]
The velocity stays finite on both sides while the shear becomes a surface measure.
Euler admits this stationary sheet because its momentum law contains no viscous term that differentiates across the jump.

Now apply the fixed-viscosity equation to this exact stationary field.
Distributionally,
\[
  \Delta u
  =
  (U_+-U_-)e_1\,\delta'_\Sigma.
\]
The stationary time derivative and transport term still vanish.
The pressure remains constant, so there is nothing on the left side of the momentum equation that can match \(\nu\Delta u\).
Allowing another pressure cannot quietly repair this particular defect: every pressure gradient has zero curl, whereas the curl of \(\nu\Delta u\) is a nonzero derivative of the sheet distribution.
The stationary Euler sheet fails the fixed-\(\nu\) Navier--Stokes momentum equation.

This gives a precise conditional class-exit witness.
If this same stationary sheet, with the same piecewise constant velocity and no compensating time dependence, is proposed as the endpoint continuation \(Q\), then
\[
  \neg\operatorname{Member}(Q)
\]
follows from the unmatched viscous distribution.
The calculation belongs to this planar slip model.
It does not prove that every nonsmooth outcome is an Euler solution, and it does not prove that every tangential discontinuity must have the same distributional structure.
The infinite planar sheet also has infinite kinetic energy on \(\mathbb R^3\), so it is not a whole-space Millennium counterexample.

What the comparison establishes is local and exact.
A bounded jump can satisfy the shared inviscid transport, pressure, and incompressibility equations while failing the fixed-viscosity equation at the interface.
The failure appears in the velocity itself, which makes it easy to see.
A finite continuous velocity can hide the same loss of smoothness in its gradient, its material acceleration, its jerk, or a still higher derivative.
Following those possibilities requires us to climb above the base value of the field.
